\section{Discussion}
\label{sec:discussion}

The results show that compact and derived open \llm artifacts have entered the public scholarly record, but not yet as well-formed curation records. A curator looking at a paper and its public traces needs to answer four practical questions: What is the artifact? Which paper relation does the repository have? What upstream relation is evidenced? What assets can be inspected or preserved? The visibility-to-curatability funnel shows that these answers often exist separately, but rarely appear together in one typed record; after usable provenance and concrete release evidence are present, the remaining drop is driven mainly by whether the release evidence reaches the asset-package level. This section translates that gap into a minimal record design.

\subsection{From Visibility to a Minimal Curatable Record}

\begin{table*}[!tbp]
\centering
\caption{A minimal curatable record for compact and derived open \llm artifacts. The record acts as a coordination contract across existing systems rather than a requirement for a new centralized registry.}
\label{tab:minimal-curatable-record}
\small
\setlength{\tabcolsep}{4pt}
\renewcommand{\arraystretch}{1.12}
\begin{tabularx}{0.98\textwidth}{@{}p{0.15\textwidth}p{0.20\textwidth}p{0.22\textwidth}p{0.20\textwidth}Y@{}}
\toprule
Field group & Core field & Primary maintainer/source & Failure prevented & Curation use \\
\midrule
Artifact identity & \texttt{artifact\_role} & Model hub / repository owner & Generic model label & Catalog type \\
\rowcolor{TableGray}
Package form & \texttt{package\_form} & Model hub / repository metadata & Hidden distribution format & Package interpretation \\
Scholarly relation & \texttt{paper\_repository\_relation} & Paper authors + scholarly index & Untyped links & Bibliographic control \\
\rowcolor{TableGray}
Upstream relation & \texttt{upstream\_relation\_type} & Model hub + paper text & Family-as-provenance & Provenance relation \\
Evidence tier & \texttt{provenance\_evidence\_tier} & Digital library / index & Overstated confidence & Confidence and scope \\
\rowcolor{TableGray}
Release assets & \texttt{release\_asset\_types} & Model hub + code repository + paper & Release-as-binary & Asset package \\
License/access signal & \texttt{license\_access\_signal} & Model hub / repository metadata & Missing governance context & Access context \\
\bottomrule
\end{tabularx}
\renewcommand{\arraystretch}{1.0}
\end{table*}


Table~\ref{tab:minimal-curatable-record} summarizes the record that follows from the measurements. It is intentionally minimal: the fields correspond to the evidence needed to move from discovery to provenance-aware curation. Artifact role and package form describe what the object is and how it is distributed. Scholarly relation links the object to papers. Upstream relation and provenance-evidence tier separate family-level discovery from stronger derivation claims. Release assets describe what can be inspected or preserved. License and access signals add governance context that should accompany, but not replace, the integrated curatability status. The added failure-prevention column makes the design logic explicit: each field blocks a specific ambiguity observed in the audit or funnel.

The table also shows that the record can be assembled as a coordination layer across existing infrastructure, provided that stable identifiers, bidirectional links, typed relation fields, release-asset fields, and versioned evidence tiers are preserved. Model hubs and repository owners are closest to artifact role, package form, upstream relation, release assets, and license metadata. Paper authors and scholarly indexes are closest to paper--repository relations. Digital libraries and preservation systems are positioned to store evidence tiers and curation-readiness status because they integrate records from multiple sources.

This design is intentionally modest. It does not require a single authority to certify every open-model artifact. Instead, it asks each infrastructure layer to expose the relation it is already closest to: model hubs expose artifact and release structure, scholarly indexes expose paper relations, and digital libraries preserve evidence tiers across systems. The result is not a heavier publication burden, but a record that makes existing public traces usable for citation, attribution, reuse assessment, and preservation triage.

The boundary cases observed in the audits explain why the fields should remain typed. A family name can mark a target model, parent, teacher, baseline, evaluator, or false positive context match. A GitHub link can expose scripts without weights. A \hf link can identify a dependency or checkpoint without an asset-level release package. These traces are useful, but they are not interchangeable. A curation record should preserve their relation type and evidentiary strength rather than compressing them into a generic ``open model'' label.

\subsection{Infrastructure Responsibilities and Design Implications}

The funnel suggests a practical division of labor. Because release assets and full recipes are rare, \textbf{model hubs} are the natural place to maintain artifact-role, package-form, upstream-relation, release-asset, and license/access fields close to the artifact. Because matched paper--\hf repository evidence covers only 5.6\% of paper records, \textbf{scholarly indexes} should type paper--repository relations rather than storing unqualified links. Because family evidence is often usable for grouping but weak for strict upstream claims, \textbf{digital libraries and preservation systems} should preserve provenance-evidence tiers and curation-readiness status across sources. The priority is not a new registry of every model artifact, but a coordination layer that keeps existing public traces joinable without losing their source or evidentiary strength.

This division of labor leads to five evidence-to-action interventions.
\begin{enumerate}[leftmargin=*,itemsep=2pt,topsep=2pt]
    \item \textbf{Role ambiguity $\rightarrow$ separate artifact role from package form.} Model hubs should distinguish base checkpoints, fine-tunes, adapters, merges, quantized files, and deployment packages without forcing them into one generic model label.
    \item \textbf{Sparse direct linkage $\rightarrow$ type paper--repository relations.} Scholarly indexes should distinguish repositories introduced by a paper, used as dependencies, evaluated as baselines, cited as background, or provided as official release packages.
    \item \textbf{Family-as-provenance risk $\rightarrow$ type upstream relations.} Model records should distinguish base, teacher, adapter target, merge input, quantization source, evaluator, baseline, and background mention relations.
    \item \textbf{Uneven evidence strength $\rightarrow$ represent provenance tiers.} Digital libraries should preserve whether a claim rests on an explicit upstream field, a matched repository with a visible upstream relation, an inventory-backed family signal, or a contextual mention.
    \item \textbf{Binary release labels $\rightarrow$ expose release packages as structured fields.} Release records should separate code, scripts, weights or checkpoints, adapters, datasets, evaluation assets, recipes, and environment details.
\end{enumerate}

These interventions align with FAIR software and research-object packaging work, which treats reuse as a property of structured packages and relations rather than of a single access flag \cite{barker2022fair4rs,soilandreyes2022rocrate}. The Model Openness Framework makes a related point at the level of openness components \cite{white2024mof}. The contribution here is empirical and record-level: open \llm artifacts are often visible, but digital-library infrastructure should make the evidence needed for citation, provenance tracking, reuse assessment, and preservation selection explicit, typed, and actionable.
